\section{Coordinate and response-slice audits}
\label{app:component-results}
Figure~\ref{fig:core_coordinate_coverage} displays all coordinatewise versus joint core events. The response-defined zero/positive slices differ from the observable covariate groups used for calibration. Paired minority improvements retain their shared-trial uncertainty.
\begin{figure}[htbp]
\centering
\includegraphics[width=\textwidth]{figures/assets/core_coordinate_coverage.pdf}
\caption{\CoreCoordinateCaption}
\label{fig:core_coordinate_coverage}
\end{figure}
\clearpage
\input{figures/tables/core_zero_slices}
\input{figures/tables/core_minority_improvements}
\clearpage
\section{Measured fitted-region cases}
\label{app:measured-regions}
These are all 32 method/case records selected before outcomes: core K 3 trial 0, test observation 0, shared across methods. Contours are approximate displays of the analyst-owned membership rule, not area integration or representative population estimates. The original recorded target's membership and center inclusion were checked for every case. Near-corner marks can overlap; exact coordinates remain in the delivered case file.
\begin{figure}[htbp]
\centering
\includegraphics[width=\textwidth]{figures/assets/regions_interior_k3.pdf}
\caption{\CoreRegionInteriorCaption}\label{fig:regions_interior}
\end{figure}
\clearpage
\begin{figure}[p]
\centering
\includegraphics[width=\textwidth]{figures/assets/regions_heterogeneous_k3.pdf}
\caption{\CoreRegionHeterogeneousCaption}\label{fig:regions_heterogeneous}
\end{figure}
\clearpage
\begin{figure}[p]
\centering
\includegraphics[width=\textwidth]{figures/assets/regions_boundary_k3.pdf}
\caption{\CoreRegionBoundaryCaption}\label{fig:regions_boundary}
\end{figure}
\clearpage
\begin{figure}[p]
\centering
\includegraphics[width=\textwidth]{figures/assets/regions_dirichlet_k3.pdf}
\caption{\CoreRegionDirichletCaption}\label{fig:regions_dirichlet}
\end{figure}
\clearpage
\begin{figure}[p]
\centering
\includegraphics[width=\textwidth]{figures/assets/regions_zeros_k3.pdf}
\caption{\CoreRegionZerosCaption}\label{fig:regions_zeros}
\end{figure}
\clearpage
\section{Numerical contrasts and sensitivity results}
\label{app:sensitivities}
The plots and tables below retain enlarged paired area uncertainty, the specific prevalence-shift comparison and the three declared sensitivity families. Cross-cell area contrasts have no within-cell numerical ranking certification. The complete machine-readable paired-contrast and allowance files preserve all comparisons, including unresolved outcomes.
\begin{figure}[htbp]
\centering
\includegraphics[width=\textwidth]{figures/assets/paired_area_resolution.pdf}
\caption{\FinalAreaCaption}
\label{fig:paired_area_resolution}
\end{figure}
\begin{figure}[htbp]
\centering
\includegraphics[width=\textwidth]{figures/assets/prevalence_shift.pdf}
\caption{\FinalShiftCaption}
\label{fig:prevalence_shift}
\end{figure}
\clearpage
\input{figures/tables/shift_coverage}
\input{figures/tables/paired_shift_effects}
\clearpage
\input{figures/tables/calibration_variability}
\input{figures/tables/small_group_cost}
\clearpage
\input{figures/tables/epsilon_sensitivity}
\clearpage
\section{Complete study tables and domain accounting}
\label{app:full-results}
All 446 declared method/group entries are retained: 426 defined and 20 explicit unsmoothed-ILR inapplicability entries, each with 100 recorded trials. Counts, actual nonempty denominators, fit/calibration counts and numerical status remain in the exact companion CSV. Zero-count groups are not assigned zero coverage; undefined positive-weight area propagates undefined. Infinite thresholds and full-domain identification bounds/unknown status are distinct diagnostics. These are descriptive pointwise intervals, not simultaneous acceptance tests.
\input{figures/tables/full/complete_overall}
\clearpage
\input{figures/tables/full/complete_subgroups}
\clearpage
\input{figures/tables/full/domain_status}
